%%%=========[EX_01]================%%%
\begin{ex}
    Phương pháp nào sau đây giúp bảo quản thực phẩm được lâu mà vẫn giữ được giá trị dinh dưỡng?
    \choice
        {\True Bảo quản trong tủ lạnh ở nhiệt độ thích hợp.}
        {Để thực phẩm ở nhiệt độ phòng trong thời gian dài.}
        {Hâm nóng lại thực phẩm nhiều lần.}
        {Để thực phẩm tiếp xúc trực tiếp với ánh nắng mặt trời.}
\end{ex}

%%%=========[EX_02]================%%%
\begin{ex}
    Nhóm thực phẩm nào sau đây chứa nhiều protein (chất đạm) nhất?
    \choice
        {Gạo, khoai, lúa mì.}
        {\True Thịt bò, trứng, cá.}
        {Dầu ăn, mỡ động vật.}
        {Cam, bưởi, rau xanh.}
\end{ex}

%%%=========[EX_03]================%%%
\begin{ex}
    Loại nhiên liệu nào sau đây được xem là năng lượng tái tạo?
    \choice
        {Than đá.}
        {Dầu mỏ.}
        {\True Năng lượng gió.}
        {Khí tự nhiên.}
\end{ex}

%%%=========[EX_04]================%%%
\begin{ex}
    Việc lạm dụng nhiên liệu hóa thạch gây ra hậu quả nào sau đây đối với môi trường?
    \choice
        {Làm giảm nhiệt độ Trái Đất.}
        {\True Tăng lượng khí thải carbon dioxide gây hiệu ứng nhà kính.}
        {Giúp không khí trở nên trong lành hơn.}
        {Làm tăng lượng oxygen trong khí quyển.}
\end{ex}

%%%=========[EX_05]================%%%
\begin{ex}
    Để bảo quản thịt tươi trong ngăn mát tủ lạnh, thời gian an toàn thông thường là bao lâu?
    \choice
        {1 -- 2 tuần.}
        {\True 3 -- 5 ngày.}
        {1 tháng.}
        {24 giờ.}
\end{ex}

%%%=========[EX_06]================%%%
\begin{ex}
    Thành phần dinh dưỡng chính của lương thực, thực phẩm cung cấp cho cơ thể con người là gì?
    \choice
        {Chất khoáng và nước.}
        {Chỉ có carbohydrate và protein.}
        {\True Carbohydrate, protein, chất béo, khoáng chất và vitamin.}
        {Chủ yếu là vitamin và tinh bột.}
\end{ex}

%%%=========[EX_07]================%%%
\begin{ex}
    Việc làm nào dưới đây là sai trong cách bảo quản lương thực, thực phẩm?
    \choice
        {Sấy khô các loại hoa quả để bảo quản lau dài.}
        {Ướp muối cho cá trước khi phơi khô.}
        {\True Để thịt ngoài không khí trong thời gian dài.}
        {Chế biến hải sản sạch sẽ rồi để trong tủ lạnh.}
\end{ex}

%%%=========[EX_08]================%%%
\begin{ex}
    Loại lương thực nào sau đây cung cấp nhiều carbohydrate (chất đường bột) nhất cho cơ thể?
    \choice
        {Thịt bò.}
        {Dầu đậu nành.}
        {\True Gạo tẻ.}
        {Cam chua.}
\end{ex}

%%%=========[EX_09]================%%%
\begin{ex}
    "Ngũ cốc" là khái niệm truyền thống thường dùng để chỉ nhóm thực vật nào?
    \choice
        {Lúa nước, ngô, dưa hấu, nhãn, vải.}
        {\True Gạo nếp, gạo tẻ, vừng, lúa mì và các loại đậu.}
        {Rau cải, bắp cải, cà chua, hành, tỏi.}
        {Khoai mỡ, sắn, củ đậu, khoai lang, bí đỏ.}
\end{ex}

%%%=========[EX_10]================%%%
\begin{ex}
    Tại Việt Nam, các loại lương thực phổ biến nhất được sử dụng trong đời sống hằng ngày là
    \choice
        {lúa mì, đại mạch, yến mạch.}
        {bánh mì, khoai lang, lúa mạch.}
        {\True lúa gạo, ngô, khoai, sắn.}
        {đậu tương, hạt dẻ, hạt điều.}
\end{ex}

%%%=========[EX_11]================%%%
\begin{ex}
    Đá cháy (băng cháy) đang được nghiên cứu như một nguồn năng lượng tiềm năng, thành phần chính của nó là
    \choice
        {chất rắn chứa oxygen dạng đông đá.}
        {hydrogen lỏng.}
        {\True methane hydrate.}
        {carbon tinh khiết.}
\end{ex}

%%%=========[EX_12]================%%%
\begin{ex}
    Để sử dụng nhiên liệu gas (đun nấu) một cách ưu việt, tiết kiệm và an toàn, chúng ta cần
    \choice
        {luôn để mức lửa lớn nhất có thể để đun nhanh.}
        {Mở lửa liên tục dù không đun gì để tiết kiệm công bật bếp.}
        {luôn luôn để ngọn lửa thật nhỏ.}
        {\True điều chỉnh mức lửa phù hợp với kích thước nồi và nhu cầu sử dụng.}
\end{ex}

%%%=========[EX_13]================%%%
\begin{ex}
    Việc thế giới bị phụ thuộc quá nhiều vào nhiên liệu hóa thạch dẫn đến rủi ro lớn nhất về mặt vĩ mô là
    \choice
        {giảm giá lương thực toàn cầu.}
        {giảm ô nhiễm môi trường do khai thác quá tải.}
        {nhiệt độ toàn cầu giảm đi do ít năng lượng hấp thụ.}
        {\True tăng rủi ro về an ninh năng lượng do cạn kiệt hoặc đứt gãy nguồn cung.}
\end{ex}

%%%=========[EX_14]================%%%
\begin{ex}
    Loại nhiên liệu nào dưới đây tạo ra ít khí thải độc hại, ô nhiễm môi trường nhất khi đốt cháy?
    \choice
        {\True Hydrogen.}
        {Than đá.}
        {Dầu mỏ.}
        {Gỗ khô.}
\end{ex}

%%%=========[BT_01]================%%%
\begin{bt}
    Một số gia đình sử dụng bếp than tổ ong để sưởi ấm trong phòng kín vào mùa đông và đã xảy ra hiện tượng ngạt khí dẫn đến tử vong. Em hãy giải thích nguyên nhân của hiện tượng này và đưa ra lời khuyên để sử dụng bếp an toàn.
    \loigiai{
        - Nguyên nhân: Khi than cháy trong phòng kín, sinh ra khí độc carbon monoxide (CO). Khí CO làm giảm khả năng vận chuyển oxygen của máu, gây ngạt thở và tử vong.
        - Lời khuyên: Tuyệt đối không đốt than, củi để sưởi ấm trong phòng kín. Cần để phòng thông thoáng hoặc sử dụng máy sưởi điện.
    }
\end{bt}

%%%=========[BT_02]================%%%
\begin{bt}
    Trong các bữa ăn tại căng tin trường học, học sinh cần làm gì để tự bảo vệ mình khỏi nguy cơ ngộ độc thực phẩm? Liệt kê ít nhất 3 biện pháp cụ thể.
    \loigiai{
        Học sinh có thể phòng tránh ngộ độc thực phẩm tại căng tin qua các biện pháp sau:
        - Rửa tay sạch bằng xà phòng trước khi ăn.
        - Quan sát kĩ màu sắc và mùi vị thức ăn; nếu có mùi lạ, ôi thiu thì tuyệt đối không ăn.
        - Chỉ mua thực phẩm đóng gói có nhãn mác rõ ràng và còn hạn sử dụng.
    }
\end{bt}
